% Proposed insertion into paper/audited_baseline.tex.
% It uses only the packages and theorem environments already present there.

\section{Mixed-rank shadows for quasi-primitive Hasse palettes}

For $m\geq1$ let $\Omega(m)=\sum_pv_p(m)$, let
$\operatorname{supp}(m)=\{p:p\mid m\}$, and, for
$E\subseteq\{2,3,\ldots\}$, put
\[
  \partial E=\{e/p:e\in E,\ p\text{ prime},\ p\mid e\}.
\]
We use the following published form of Katona's intersecting-shadow theorem:
if $\mathcal F$ is a finite intersecting family of $s$-subsets, then
\[
  \card{\partial_{\rm set}\mathcal F}\geq\card{\mathcal F}.
  \tag{KS}
\]
See \cite{Katona1964}; the more general statement in
\cite[Theorem~1.3]{FranklKatona2021}, with $t=1$ and $\ell=s-1$, also gives
(KS).  The case $s=1$ is immediate.  Notice that (KS) applies to every
subfamily of an intersecting family.

\begin{theorem}[Mixed-rank two-shadow map]\label{thm:mixed-two-shadow}
Let $E\subseteq\{2,3,\ldots\}$ be finite or countable and satisfy
$\gcd(e,f)>1$ for all distinct $e,f\in E$.  There are maps
\[
  \psi:E\longrightarrow\N,
  \qquad \tau:E\longrightarrow\{0,1\},
\]
such that $e/\psi(e)$ is prime and
\[
  e\longmapsto(\psi(e),\tau(e))
  \tag{MS}
\]
is injective.  Equivalently, one may choose an immediate multiplicative
shadow of each member so that every integer is chosen at most twice.  Hence
$\card G\leq2\card{\partial G}$ for every finite $G\subseteq E$.
\end{theorem}

\begin{proof}
First suppose that $E$ is finite and all its members have the same value
$\Omega(e)=k$.  Fix a support size $s$ and set
\[
  \mathcal A_s=\{\operatorname{supp}(e):e\in E,
                  \ \card{\operatorname{supp}(e)}=s\}.
\]
This set family is intersecting.  By (KS), every
$\mathcal G\subseteq\mathcal A_s$ satisfies
$\card{\partial_{\rm set}\mathcal G}\geq\card{\mathcal G}$; Hall's
theorem therefore gives an injection
\[
  T_s:\mathcal A_s\longrightarrow\partial_{\rm set}\mathcal A_s,
  \qquad T_s(S)\subset S,
  \qquad\card{S\setminus T_s(S)}=1.
\]
Let $q_s(S)$ be the unique prime in $S\setminus T_s(S)$ and define
$\psi_s(e)=e/q_s(S)$ when $\operatorname{supp}(e)=S$.

The map $\psi_s$ is injective on this support-size slice.  Indeed, if the
deleted prime occurs in $e$ with exponent at least two, the support of the
image is $S$; if it occurs with exponent one, the support of the image is
$T_s(S)$.  Two images of the first kind determine the same $S$, and two of
the second kind determine the same $S$ by injectivity of $T_s$; in either
case the deleted prime and then $e$ are determined.  Images of the two kinds
cannot coincide because their support sizes are $s$ and $s-1$.

Now fix an output $d$ with support size $r$.  It can have a preimage only
from the support-size $r$ slice, when a repeated prime was deleted, or from
the support-size $r+1$ slice, when a prime disappeared.  Each slice supplies
at most one preimage, so the combined map has fibres of size at most two.
For a finite mixed-rank family, apply this construction separately to
$E_k=\{e\in E:\Omega(e)=k\}$.  Its images all have rank $k-1$, so images
coming from different $k$ do not meet.  Colour the at most two members of
each fibre by $0$ and $1$ to obtain (MS).

Finally let $E$ be countable.  Make a bipartite graph with left side $E$,
right side $\N\times\{0,1\}$, and an edge
$e\sim(d,t)$ exactly when $e/d$ is prime.  Every left degree is finite.
The finite case verifies Hall's condition for every finite left subset.
Matching each finite initial segment and applying K\"onig's infinity lemma
to the resulting finitely branching tree gives a matching of all of $E$,
which is (MS).  Restriction to a finite $G$ gives
$\card G\leq2\card{\partial G}$.
\end{proof}

This theorem applies canonically to the local Hasse palettes of a
quasi-primitive set.  In the divisibility poset induced by $B$, let
$L_B(b)$ denote the lower covers of $b$.  Every nonminimal $b$ has such a
cover because $b$ has only finitely many divisors.  Choose
$\sigma(b)\in L_B(b)$ for every nonminimal $b$, let $R$ be the minimal
members of $B$, and put
\[
  E_d=\{b/d:\sigma(b)=d\}.
\]

\begin{corollary}[Hasse palette compression]\label{cor:hasse-palette}
If $B$ is finite or countable and quasi-primitive, every $E_d$ is primitive
and pairwise non-coprime.  Moreover, there are maps
$s:B\setminus R\to\N$ and $t:B\setminus R\to\{0,1\}$ such that
\[
  b\longmapsto(\sigma(b),s(b),t(b))
  \quad\text{is injective},
  \qquad
  \sigma(b)s(b)=\frac{b}{p_b}\leq\frac b2
  \tag{HC}
\]
for a prime $p_b\mid b/\sigma(b)$.
\end{corollary}

\begin{proof}
If distinct $e_1,e_2\in E_d$ satisfied $e_1\mid e_2$, then
$d<de_1<de_2$ and $d\mid de_1\mid de_2$, contrary to $d$ being a lower
cover of $de_2$.  Thus $E_d$ is primitive.  If
$\gcd(e_1,e_2)=1$, the preceding conclusion makes $de_1,de_2$
incomparable, while
$\gcd(de_1,de_2)=d\in B$, contradicting quasi-primitivity.  Apply
Theorem~\ref{thm:mixed-two-shadow} to each $E_d$, and for $b=de$ set
$s(b)=\psi_d(e)$ and $t(b)=\tau_d(e)$.  The parent label separates the
palettes, (MS) gives injectivity within each palette, and deleting one prime
gives (HC).
\end{proof}

The binary type in this statement is not a proof artefact.

\begin{proposition}[Strict one-type obstruction]\label{prop:one-type-barrier}
For distinct primes $p,q,r$, the family
\[
 E^\ast=\{qr^2,q^2r,pr^2,pqr,pq^2,p^2r,p^2q\}
\]
is primitive and pairwise non-coprime, but it has no injective map
$\varphi:E^\ast\to\partial E^\ast$ with $e/\varphi(e)$ prime.
\end{proposition}

\begin{proof}
All seven members have $\Omega=3$, so none divides another.  Their prime
supports are the three two-element subsets of $\{p,q,r\}$, together with
$\{p,q,r\}$, and hence meet pairwise.  Directly,
\[
  \partial E^\ast=\{p^2,q^2,r^2,pq,pr,qr\}.
\]
Thus an immediate-shadow injection would map seven objects into six.
\end{proof}

Proposition~\ref{prop:one-type-barrier} excludes every universal one-type
immediate-shadow scheme; it does not show that the coefficient $2$ in the
finite inequality is the best possible real constant.  More importantly,
Corollary~\ref{cor:hasse-palette} is still only an \emph{abstract labelled}
compression.  Shadow values from different ranks can be comparable (for
example, the valid choices $6\mapsto2$ and $20\mapsto4$), and labels attached
to different parents can collide.  The unproved step is a bounded-cost
arithmetic erasure of the parent and type labels that leaves a primitive
ordered sequence.  Consequently these results do not settle the
quasi-primitive case of Problem~892.

\section{Moving-endpoint avoidance at the logarithmic barrier}

Consider the test profile
\[
 n_j=\left\lceil j+j(\log(j+1))^2\right\rceil,
 \qquad X_j=2^{n_j}.
\]
For $i<k$, an odd $d\in(X_i/2,X_i]$, and $1\leq h\leq i$, put
\[
 Y=Y_{i,k,d}=X_k/d,
 \qquad
 D_{i,h}(d)=\left(X_{i+h}/(2d),X_{i+h}/d\right]
                 \cap\N_{\rm odd},
\]
and define the unconditioned finite-block avoidance set
\[
 E^\square_{i,k}(d)=
 \left\{r\in(Y/2,Y]\cap\N_{\rm odd}:
 a\nmid r\ \text{for all }a\in\bigcup_{h\leq i}D_{i,h}(d)\right\}.
 \tag{A}
\]
For a middle anchor $i<k/2$, this is the quotient set appearing in the
maximal-old-divisor reduction of the stage manuscript, but (A) omits both
maximal-anchor survival and terminal-pool rank conditions.

The next result uses precisely these standard analytic inputs, uniformly as
their arguments tend to infinity:
\[
 \sum_{p\leq x}\frac1p=\log\log x+O(1),\qquad
 \sum_{x<p\leq2x}\frac1p\ll\frac1{\log(2x)},
 \tag{AI1}
\]
\[
 \sum_{\substack{x<u\leq2x\\\Omega(u)=2}}\frac1u
 \ll\frac{\log\log(3x)}{\log(2x)},
 \qquad
 \pi(x)-\pi(x/2)\gg\frac{x}{\log x}.
 \tag{AI2}
\]
Here the first formula is Mertens' theorem.  The second is the standard
dyadic prime-reciprocal estimate (for example, a consequence of the prime
number theorem), the first estimate in (AI2) is the standard dyadic
semiprime upper bound, and the last is the prime number theorem in dyadic
intervals.  For completeness, the semiprime bound follows by writing
$u=pq$ with $p\leq q$, applying the usual PNT upper bound to $q$ for each
$p\leq\sqrt{2x}$, and then using Mertens' estimate.  No prime theorem in an
interval shorter than a fixed positive proportion of its endpoint is used.

\begin{proposition}[Moving-endpoint semiprime floor, modulo (AI1)--(AI2)]
\label{prop:semiprime-floor}
Assume the named standard estimates (AI1)--(AI2).  Fix $0<a<b<1/2$.
There is $c_{a,b}>0$ such that, for all sufficiently
large $k$, all integers $i$ with $ak\leq i\leq bk$, and every odd
$d\in(X_i/2,X_i]$,
\[
  \card{E^\square_{i,k}(d)}\geq c_{a,b}\frac{Y_{i,k,d}}i.
  \tag{SF}
\]
\end{proposition}

\begin{proof}
Write $Z=X_{2i}/d$ and
$\mathcal D=\bigcup_{h\leq i}D_{i,h}(d)$.  If
$\Theta_h=X_{i+h}/(2d)$, then
$D_{i,h}(d)=(\Theta_h,2\Theta_h]\cap\N_{\rm odd}$ and
$\mathcal D\subset(1,Z]$.  The mean-value theorem applied to
$x+x(\log(x+1))^2$, with the ceiling errors absorbed, gives uniformly for
$1\leq h\leq i$
\[
  \log\Theta_h\asymp h(\log i)^2,
  \qquad \log Z\asymp i(\log i)^2.
  \tag{1}
\]
Regular variation, uniformly for $ak\leq i\leq bk$, gives
\[
 \frac{\log Y}{\log Z}
 =\frac{n_k-n_i+O(1)}{n_{2i}-n_i+O(1)}
 =\frac{k-i}{i}\bigl(1+o_{a,b}(1)\bigr).
\]
Since $(k-i)/i\geq(1-b)/b>1$, there is a constant
$\delta_{a,b}>0$ such that, for all sufficiently large $k$,
\[
  \frac{\log Y}{\log Z}\geq1+\delta_{a,b}.
  \tag{2}
\]

Choose $0<\eta<\min\{1,\delta_{a,b}/2\}$ and put
$P=Z^{\eta/2}$.  Let $\mathcal V$ be the products $u=pq$ of two distinct
odd primes $p,q\leq P$ for which none of $p,q,pq$ belongs to $\mathcal D$.
By (AI1),
\[
 \sum_{\substack{p\leq P\\p\ \text{odd prime}}}\frac1p\asymp\log i.
 \tag{3}
\]
On the other hand, (1), (AI1), and (AI2) imply
\[
 \sum_{h\leq i}\ \sum_{\Theta_h<p\leq2\Theta_h}\frac1p
 \ll\frac1{\log i},
 \tag{4}
\]
and
\[
 \sum_{h\leq i}\ \sum_{\substack{\Theta_h<u\leq2\Theta_h\\
                                    \Omega(u)=2}}\frac1u
 \ll\frac1{(\log i)^2}
       \sum_{h\leq i}\frac{\log(3h(\log i)^2)}h
 \ll1.
 \tag{5}
\]
The reciprocal weight of all products of two distinct odd primes at most
$P$ is $\asymp(\log i)^2$.  By (4), products for which one prime lies in
$\mathcal D$ have total weight $O(1)$; by (5), the same is true for
products which themselves lie in $\mathcal D$.  Therefore
\[
  \sum_{u\in\mathcal V}\frac1u\gg(\log i)^2.
  \tag{6}
\]

For $u\in\mathcal V$ we have $u\leq Z^\eta$, so (2) gives
$x_u:=Y/u\geq Z^{1+\delta_{a,b}/2}$.  In particular $x_u/2>Z$ for large
$k$.  The dyadic prime number theorem in (AI2), together with
$\log x_u\asymp_{a,b}\log Z$, supplies
\[
 \#\{s:Y/(2u)<s\leq Y/u,\ s\text{ prime}\}
 \gg_{a,b}\frac{Y}{u\log Z}.
 \tag{7}
\]
Every such $s$ exceeds $Z$.  Hence in $r=us$ the factor $u$ is exactly the
part supported on primes at most $Z$, so representations belonging to
different $u$ are distinct.  Any divisor of $r$ not exceeding $Z$ divides
$u=pq$; its possibilities are $1,p,q,pq$, none of which lies in
$\mathcal D$.  Thus every $r$ counted in (7) belongs to (A).  Summing (7)
over $u$, and using (1) and (6), proves
\[
 \card{E^\square_{i,k}(d)}
 \gg_{a,b}\frac{Y}{\log Z}(\log i)^2
 \gg_{a,b}\frac Yi.
\]
\end{proof}

It follows that the avoidance density for each middle anchor is at least
$c_{a,b}/i$; consequently its sum over $ak\leq i\leq bk$ stays bounded
away from zero.  Thus neither a supremum-over-$d$ anchorwise sum nor its
unconditioned mean-over-$d$ analogue can yield $o(1)$.  This is a strict
obstruction to that proof mechanism, not evidence that the profile is
inadmissible: a terminal integer may be counted by many anchors, and the
maximal-anchor and rank conditions omitted from (A) can change the actual
union.

We finally record a transfer lemma that does not require a full period of
the least common multiple.  For $Y\geq3$ let
$I_Y=(Y/2,Y]\cap\N_{\rm odd}$ and $N_Y=\card{I_Y}$.  If $F$ is a finite
set of odd integers greater than one and $s\geq1$, define
\[
 \mathcal B_{2s}(F)=1+
 \sum_{\substack{\varnothing\ne S\subseteq F\\\card S\leq2s}}
 (-1)^{\card S}\frac1{\operatorname{lcm}(S)}.
\]

\begin{lemma}[Short-interval Bonferroni transfer]
\label{lem:bonferroni-transfer}
Uniformly in $Y,F,s$,
\[
 \frac1{N_Y}\card{\{r\in I_Y:f\nmid r\text{ for every }f\in F\}}
 \leq \mathcal B_{2s}(F)
 +O\!\left(\frac1Y\sum_{j=0}^{2s}\binom{\card F}{j}\right).
 \tag{BF}
\]
If $2s\geq\card F$, inclusion--exclusion with the exact interval counts is
an identity.
\end{lemma}

\begin{proof}
For every odd $q\geq1$, direct counting gives, uniformly also for $q>Y$,
\[
 \frac1{N_Y}\card{\{r\in I_Y:q\mid r\}}=\frac1q+O(1/Y).
 \tag{8}
\]
Apply the even Bonferroni upper bound to the complement of the union of the
events $\{f\mid r\}$, $f\in F$.  The intersection indexed by
$S\subseteq F$ is $\{\operatorname{lcm}(S)\mid r\}$.  Substitution of
(8) in the retained terms yields (BF), since their number is
$\sum_{j\leq2s}\binom{\card F}{j}$.  Keeping every term gives the stated
identity.
\end{proof}

For a finite odd generator set $F$, write
$E(F)=\{r\in\N_{\rm odd}:f\nmid r\text{ for every }f\in F\}$.

\begin{remark}[Conditional actual-union proof interface]
The stage manuscript proposes a finite-block reduction with adjacent- and
terminal-strip errors $o(X_k)$.  That reduction has not been admitted as a
theorem here.  To state its remaining target without a free index, let
$\mathcal R_k$ denote the finite set of pairs $(i,d)$ retained by its
maximal-anchor bookkeeping, where $i<k/2$ and $d\in(X_i/2,X_i]$ is odd.
For each $(i,d)\in\mathcal R_k$, choose a finite
\[
 F_{i,k,d}\subseteq\bigcup_{1\leq h\leq i}D_{i,h}(d).
\]
The missing sufficient estimate inside that proposed reduction is
\[
 \lim_{J\to\infty}\sup_{k>J}\frac1{X_k}
 \left|\bigcup_{\substack{(i,d)\in\mathcal R_k\\i\geq J}}
 d\bigl(I_{Y_{i,k,d}}\cap E(F_{i,k,d})\bigr)\right|=0.
 \tag{AU}
\]
An actual quotient avoids the full generator union and hence every selected
subset $F_{i,k,d}$, so (AU), together with the two asserted $o(X_k)$ error
estimates, is exactly the terminal estimate required by that conditional
argument.  This paragraph records a proof interface; it neither proves (AU)
nor promotes the unaudited finite-block reduction to a result of this paper.
\end{remark}

Lemma~\ref{lem:bonferroni-transfer} only transfers a finite certificate to
the moving interval; it neither makes $\mathcal B_{2s}(F)$ small nor proves
the actual-union estimate (AU).
Proposition~\ref{prop:semiprime-floor} rules out replacing that estimate by
a sum of unconditioned anchorwise bounds.  Thus neither admissibility nor
inadmissibility of the logarithmic-barrier profile, and no solution of
Problem~892, follows here.
